\documentclass[11pt]{article}
\usepackage{amsmath,amssymb,amsthm,geometry,hyperref,booktabs}
\geometry{margin=1in}
\newtheorem{theorem}{Theorem}
\newtheorem{proposition}{Proposition}
\newcommand{\C}{\mathbb C}
\newcommand{\OO}{\mathcal O}
\newcommand{\crit}{\mathfrak c}
\newcommand{\MF}{\mathcal F}
\newcommand{\ord}{\operatorname{ord}}
\title{The prepared-suspension formula and the published\\
composition and polar-intersection theorems}
\author{Campaign 9000059: primary-source comparison}
\date{November 5, 2026}
\begin{document}
\maketitle

\begin{abstract}
The already reviewed claim C059LOCAL gives a uniform local
Milnor-number formula for a prepared analytic suspension.  We identify
its precise relationship to N\'emethi's ICIS composition theorem:
the candidate evaluates a residual polar intersection which that
theorem leaves as an input.  We also compute the removed critical-cycle
contribution and explain exactly why the applicable Iomdin--L\^e formula
does not determine the remaining intersection.  Independent-block
Sebastiani--Thom statements and generic-hyperplane formulas cannot be
applied by simply ignoring the constraint or its hypotheses.

This identifies a mathematically substantive \emph{candidate}
contribution, not established originality or publishability.  The
comparison below neither proves direct subsumption by a stronger
published numerical evaluation nor establishes that no such evaluation
exists.  The publication gate remains open.  The accepted local theorem,
its proof and certificate, and the reviewed M45 verification are
preserved.
\end{abstract}

\section{Exact candidate and the retained evidence}

All rings below are convergent local analytic rings at the origin.
The Milnor number of an isolated hypersurface germ is the dimension
of its Jacobian quotient; its equality with the bouquet rank is the
previously checked isolated-Milnor-fibre fact in
\cite[Section 1, items (4)--(5), p.~6]{massey}.

\begin{theorem}[Retained C059LOCAL v2]\label{thm:candidate}
Let \(p=(v,u)\), and let \(h:(\C^2,0)\to(\C,0)\) have an isolated
critical point.  Suppose that every nonconstant analytic \(p\)-arc
satisfies
\[
 \nu(\nabla h)\leq L\min(\nu v,\nu u).
\]
Let \(k\geq1\) and
\[
 R(p,w)=w^k+\sum_{j<k}A_j(p)w^j,\qquad A_j(0)=0,\qquad R_v(0)\ne0.
\]
The surface \(R=0\) is smooth, with coordinates \((u,w)\).
Assume its restriction \(d=h|_{\{R=0\}}\) has an isolated critical
point and \(d(u,0)\not\equiv0\).
For integers \(a,b>1\) with
\[
 \frac{ab-a-b}{a+b}>L,
\]
the function
\[
 f=\left.h(p)+w(y^a+z^b)\right|_{X},\qquad
 X=\{yz=R(p,w)\}\subset(\C^5,0),
\]
has an isolated critical point, and
\begin{equation}\label{eq:candidate}
 \mu(f)=\mu(d)+(kab+a+b)\mu(h).
\end{equation}
\end{theorem}

The proof is the general intersection-length theorem in
\texttt{research/audit\_report.tex}, Sections 2--3.  Its arc and
finite-flatness argument works with the exact analytic \(R\).
The supporting \texttt{research/verify.py} and
\texttt{research/exact\_certificate.json} certify the specified algebraic
inputs and identities, not the analytic proof in a proof kernel.
The Engineer and independent Reviewer judgements on C059LOCAL v2
are retained without alteration.  No unchanged certificate replay is
presented as a new mathematical result here.

For comparison with the staged manuscript, its M45 Section 04
uses the particular base \(T(X-T^2)(X-2T^2)\), two degree-four
prepared graph maps, their particular plane restrictions, and a
small deformation parameter.  Equation \eqref{eq:candidate}
allows arbitrary \(h,k,R\) satisfying the displayed hypotheses and
does not require a tube or lattice identification.  That is a genuine
scope difference from this manuscript's special case.  It does
\emph{not}, by itself, establish a difference from the published
singularity-theory literature.

\section{What the primary theorems actually say}

\subsection{Local and global independent-block composition}

N\'emethi \cite[Theorem L1 and Corollary 2.3, pp.~2--3]{nemethi80}
takes arbitrary analytic germs \(g(x)\) and \(j(y)\) in
\emph{independent variable blocks}, and an analytic plane germ
\(P(c,e)\).  Write \(G,J,Q_P\) for their Milnor fibres.
If \(n_c=0\) when \(c\mid P\), and otherwise
\(n_c=I_0(P,c)\), with \(n_e\) defined similarly, the formula is
\begin{equation}\label{eq:local-ts}
 \chi(\MF_{P(g,j)})=
 (\chi(Q_P)-n_c-n_e)\chi(G)\chi(J)
       +n_c\chi(J)+n_e\chi(G).
\end{equation}
The initial connectedness assumption on \(Q_P\) is removed
componentwise in Remark 2.2.  Isolatedness of \(g,j\) is not required.
For \(P=c+e\) this gives the join; for \(P=ce\), the fibre is a
bundle over an annulus and has Euler characteristic zero.
For \emph{isolated} sums the familiar numerical consequence is
\(\mu(g+j)=\mu(g)\mu(j)\), not an additive formula.
It is not legitimate to use a finite isolated Milnor number for
the nonisolated factor \(w(y^a+z^b)\).

The global sum theorem \cite[Theorems 2.2 and 2.4]{nemethi-global}
concerns polynomials on \(\C^n\times\C^m\) and their generic
\emph{global} fibres, outside the sum of the bifurcation sets.
These sets include atypical values, not just critical values.
There is no global properness hypothesis in that theorem.
The generalized global result
\cite[Section 3.2 and Theorem G1, p.~6]{nemethi80}
also uses independent-block polynomials.  Its plane polynomial
depends effectively on both variables, its generic fibre is
connected, and it requires
\[
 T_c\cap\Lambda_g=\varnothing,\qquad
 T_e\cap\Lambda_j=\varnothing,
\]
where \(T_c,T_e\) specify the tangent lines at the two indicated
points at infinity of the projective generic plane fibre.
These statements do not identify a local analytic prepared germ
with a global polynomial fibre.

The constraint \(yz=R(p,w)\) couples the proposed blocks.  Solving
it for \(v\) gives
\[
 h(v(u,w,yz),u)+w(y^a+z^b),
\]
not a displayed independent-block sum.
This rules out that \emph{direct substitution};
it is not a theorem that no indirect reduction is possible.

\subsection{The strongest applicable composition statement}

N\'emethi \cite[Sections 2.0.1--2.0.4 and Theorem A(a),
pp.~66--69]{nemethi79} allows shared variables, but assumes that
\[
 U=(H,G):(\C^{n+1},0)\longrightarrow(\C^2,0)
\]
defines an isolated complete intersection singularity (ICIS).
Let \(\crit=I_2 DU\), and put
\(\gamma_K(\crit)=\crit:K^\infty\) for \(K=P(H,G)\).
The theorem gives
\begin{equation}\label{eq:nemethi}
 \chi(\MF_K)=\chi(\MF_P)\chi(\MF_U)
               +(-1)^n I_0(\gamma_K(\crit),K).
\end{equation}
The unsaturated quotient \(\OO/\crit\) is a
one-dimensional Cohen--Macaulay ring; the saturation removes
the critical-curve components contained in \(K=0\).
Part (a) does not assume these components are absent.
The stronger bouquet assertion in part (b) does.  In particular,
``no common discriminant component'' is not a condition of
part (a), nor should it be replaced by transversality at the origin.

\begin{proposition}[An exact reduction, not a novelty certificate]
\label{prop:reduction}
Under Theorem~\ref{thm:candidate}, define on the unconstrained
ambient germ \(\C^5\)
\[
 Q=y^a+z^b,\qquad F=h(v,u)+wQ,\qquad S=yz-R(v,u,w).
\]
Let \(\crit=I_2D(F,S)\) and let \(\Gamma\) be the residual critical
cycle defined by \(\crit:F^\infty\).  Then
\begin{equation}\label{eq:residual-F}
 \mu(f)=\mu(h)+I_0(\Gamma,F).
\end{equation}
Consequently the retained formula is equivalent to the evaluation
\begin{equation}\label{eq:delta-F}
 I_0(\Gamma,F)=\mu(d)+(kab+a+b-1)\mu(h).
\end{equation}
\end{proposition}

\begin{proof}
Apply \eqref{eq:local-ts} first to the product \(wQ\):
\(\chi(\MF_{wQ})=0\).  Then apply its sum case to the
independent blocks \(p\) and \((w,y,z)\).  It gives
\(\chi(\MF_F)=1-\mu(h)\).
Since \(S_v=-R_v\) is a unit, \(S=0\) is smooth.
The accepted isolatedness of \(f=F|_{S=0}\) makes
\((F,S)\) an ICIS.  Its Milnor fibre is that of \(f\), and has
Euler characteristic \(1-\mu(f)\).
Use \eqref{eq:nemethi} with \(n=4\) and \(P(c,e)=c\), whose
Milnor fibre is a disk.  Thus
\[
 1-\mu(h)=1-\mu(f)+I_0(\Gamma,F).
\]
This proves \eqref{eq:residual-F}; substitution of the retained
theorem proves \eqref{eq:delta-F}.
\end{proof}

In particular, this application uses the already accepted
isolatedness.  It is not a replacement proof of isolatedness.
All prepared coefficients are retained: the unit \(R_v\) shows
that the exact minor ideal is generated by
\begin{equation}\label{eq:minors}
\begin{split}
 R_vh_u-R_uh_v,\quad &R_vQ-R_wh_v,\\
 aR_vwy^{a-1}+zh_v,\quad &bR_vwz^{b-1}+yh_v.
\end{split}
\end{equation}
The remaining minors are combinations of these after division
by \(R_v\).  A bounded exact Jacobian adapter calculation verifies
the sign \(R_vQ-R_wh_v\); its typed contract, payload, output,
server version, and execution status are retained in
\path{research/literature/prepared_comparison_jacobian_*.json}.
This determinant calculation is not evidence of originality or
a mechanical certification of Proposition~\ref{prop:reduction}.

\section{Iomdin--L\^e and the precise remaining polar invariant}

A tempting application on \(X\) would use the pair
\((h(p),wQ)\).  It is \emph{not} an ICIS: its zero fibre contains
the two curves
\[
 p=w=z=0,\qquad p=w=y=0.
\]
At every point on either curve \(d(h(p))=0\), so the pair has
rank less than two away from the origin.  This defeats the
direct use of both the power-series construction
\cite[Section 2.1.4]{nemethi79} and the general high-order
perturbation variant in its Section 2.1.7.  The latter
explicitly retains the ICIS hypothesis, requires the perturbed
function to be isolated, and imposes an order bound involving
the residual polar length.  High exponents alone do not repair
the failed hypothesis.

There is, however, a valid Iomdin--L\^e application to the ambient
pair \((F,S)\).  Computing exactly what it says gives a stronger
comparison than merely rejecting the invalid pair.

\begin{proposition}[The removed cycle and ambient Iomdin formula]
\label{prop:iomdin}
Let
\[
 B=\bigl(k(a-1)(b-1)+a+b\bigr)\mu(h).
\]
For the pair \((F,S)\) above, the critical components removed by
saturation have total intersection \(B\) with \(S=0\), and
\begin{align}
 \mu(f)&=I_0(\Gamma,S)+B,\label{eq:full-polar}\\
 \mu(F+S^q)&=qB-\mu(h)\qquad(q\ \hbox{sufficiently large}).
 \label{eq:iomdin}
\end{align}
Thus the precise additional numerical evaluation supplied by the
retained theorem is equivalently
\begin{equation}\label{eq:delta-S}
 \boxed{I_0(\Gamma,S)=\mu(d)+k(a+b-1)\mu(h).}
\end{equation}
\end{proposition}

\begin{proof}
Direct differentiation gives
\[
 \Sigma F=\{p=y=z=0\}\ \cup\
              \{p=w=0,\ y^a+z^b=0\}.
\]
The first component is the \(w\)-axis.  At its nonzero points,
the transversal Milnor number is
\((a-1)(b-1)\mu(h)\).  Indeed \(w\) is a unit there;
on a transverse \(S\)-slice, take analytic roots of this unit
in the changes \(Y=yw^{1/a}\), \(Z=zw^{1/b}\).
The transversal germ becomes \(h(p)+Y^a+Z^b\).
The order of \(S\) on this axis is \(k\), since
\(R(0,w)=w^k\).

At a nonzero point of a branch of the second component, \(Q\)
is a smooth coordinate and the transversal germ is \(h(p)+wQ\).
Its transversal Milnor number is \(\mu(h)\).
If \(g=\gcd(a,b)\), the reduced curve \(Q=0\) has \(g\)
branches, with parametrizations of orders \(b/g,a/g\)
in \(y,z\).  On each, \(S=yz\) has order \((a+b)/g\).
Thus these branches contribute \((a+b)\mu(h)\).
The restriction of \(S\) to each critical branch is nonconstant
and is a local coordinate away from the origin.

There are no additional critical-curve components contained in
\(F=0\).  On any such component, the ICIS condition makes \(S\)
nonconstant.  Generically \(dF=\lambda\,dS\); differentiation
along the component, on which \(F\) is constant, gives
\(\lambda=0\).  It is therefore a component of \(\Sigma F\).
The generic scheme multiplicities of \(\crit\) on these
components are exactly the transversal Jacobian lengths just
computed.  This proves the claimed value \(B\).

Using \(S\) as one analytic coordinate, the ideal
\((S,\crit)\) becomes the Jacobian ideal of \(f\) together with
that coordinate.  Its length is \(\mu(f)\).  Additivity of
intersection numbers for the one-dimensional critical cycle,
with its scheme multiplicities, gives \eqref{eq:full-polar}.
In particular neither a reduced support nor an unsaturated
length is being substituted for the residual cycle.

Section 2.1.4 of \cite{nemethi79}, applied to this ICIS pair,
gives for sufficiently large \(q\)
\[
 \mu(F+S^q)+\mu(f)=I_0(\Gamma,F)+qB.
\]
Equation \eqref{eq:residual-F} cancels \(\mu(f)\) from this
identity, proving \eqref{eq:iomdin}.  Finally,
\[
 kab+a+b-\bigl(k(a-1)(b-1)+a+b\bigr)=k(a+b-1)
\]
and \eqref{eq:candidate} yield \eqref{eq:delta-S}.
\end{proof}

The applicable high-power formula \eqref{eq:iomdin} does not even
contain \(\mu(d)\).  Combining these identities alone cannot
recover \eqref{eq:delta-S}, because the unknown section Milnor
number cancels.  This is a limited algebraic conclusion about
these formulas, not an exclusion of other uses of polar or
Iomdin--L\^e theory.

\section{Generic sections, L\^e--Greuel, and the accepted calculation}

L\^e's original attaching result
\cite[Theorem (2.3) and Section (3.1), pp.~267--268]{le}
constructs the Milnor fibre from a generic hyperplane-section
fibre by attaching \(I_0(\Gamma,F)\) cells.  Its choice of
hyperplane uses Theorem (1.1) and the good-stratification
transversality conditions in Section 2.  Merely knowing that
\(S_v\) is a unit does not assert those genericity conditions
for this particular nonlinear \(S\).
The ICIS composition theorem above is the applicable
shared-variable replacement.  Likewise, the Teissier identity
in \cite[Exercise 3.3, p.~19]{massey} is
\[
 I_0(\Gamma_{G,z_0},G)=I_0(\Gamma_{G,z_0},z_0)
                  +I_0(\Gamma_{G,z_0},G_{z_0}),
\]
when the indicated polar intersections are proper.  It
relates polar intersections, rather than evaluating an arbitrary
one from \(k,\mu(h),\mu(d)\).

A modern, directly applicable L\^e--Greuel formulation is
Damon--Pike \cite[Theorem 3.3, equation (3-2), p.~923]{damon-pike}.
It assumes an \(H\)-holonomic free divisor in \(\C^N\), an
ICIS \(X=\varphi^{-1}(0)\subset\C^n\), and transversality of the
map \(g|_X\) to that divisor away from the origin.
Put
\[
 n=5,\quad p=N=1,\quad \varphi=S,\quad g=F,\quad
 H(t)=t,\quad V=\{0\}.
\]
Here \(V\) is \(H\)-holonomic, its annihilating vector-field
module is zero, \(X\) is smooth with \(\mu(S)=0\), and
isolatedness of \(f\) supplies the required transversality.
The formula specializes to
\[
 \mu(f)=
 \dim_\C
 \frac{\OO_{X,0}^2}
 {\langle(S_{x_i},F_{x_i}):1\leq i\leq5\rangle}
 =
 \dim_\C\frac{\OO_{\C^5,0}}{(S,\crit)}.
\]
The last equality follows by eliminating the module coordinate
with the unit \(S_v\).  It is also immediate by the analytic
chain rule, independently of this general theorem.
Again the numerical evaluation, rather than the description
of the invariant as a length, remains to be supplied.

The accepted proof supplies that evaluation without computing
the discriminant or discarding a Weierstrass unit.  For a
precise reference to its independent mathematical work, write
\(m=a+b\), \(n=ab\), \(N=n-m\), and
\(\ell=\ord d(u,0)-1\).  Its three exact local lengths are
\[
 L_z=a\ell+(a-1)\mu(d),\qquad
 L_w=a\ell+m\mu(h),\qquad
 L_A=a\mu(d)+kn\mu(h).
\]
It proves finiteness before using
\(\mu(f)+L_z=L_w+L_A\).
The nontrivial uniform analytic input is the polar-intersection
lemma in Section 2 of the retained proof: the arc-gradient
bound prevents cancellation under the perturbation, and
finite flatness preserves length.  The finite free cover
\(y=Y^b,z=Z^a\) then makes the last length accessible.
These are references to the accepted proof, not another proof
or another computational certificate of the same theorem.

\section{Contribution and publication boundary}

The strongest applicable published result inspected is
\eqref{eq:nemethi}.  It is much more general than the candidate
as a structural theorem: arbitrary ICIS pairs and analytic
plane compositions are allowed.  Calling C059LOCAL a new
generalization of that theory would be wrong.
The concrete possible contribution is instead the
\emph{explicit evaluation} \eqref{eq:delta-S}, for the
specified prepared family and explicit sufficient gradient
condition.  It replaces a residual polar intersection in five
variables by two plane Jacobian lengths and \(k,a,b\), for
analytic rather than merely polynomial coefficients.
Together with the independent local proof this is more than
matching examples.  However, \eqref{eq:delta-S} is
\emph{equivalent to the retained theorem}; changing its notation
does not create another original result.

Relative to the strongest prior \emph{campaign} result, there
is no stronger correctness theorem here.  The improvement is
the exact comparison: which published theorem applies, which
ideal and sign it uses, the value of the removed cycle, and
the residual numerical input that remains.  All independently
reviewed M45 conclusions remain as before.

Technical substance is therefore identifiable, but
field-level significance is not yet demonstrated.  A useful
uniform computation need not be a publishable original
contribution.  No sharpness result, new topological
classification, new integral lattice theorem, or independent
application beyond the retained construction has been proved
in this mission.  None is silently added to rescue the
publication claim.

The literature comparison was bounded.  It started with both
N\'emethi references in the retained Massey text and the
specified Compositio 79 paper, then inspected the original
L\^e statement and the directly relevant Damon--Pike
L\^e--Greuel formula.  A separately delegated nonlinear-section
lead found a discriminant freeness theorem, not a numerical
polar evaluation; it is not used as a proof input.
No directly subsuming \emph{numerical} theorem was verified in
this examination.  This is not a claim that none exists.
The preliminary search synthesis had no usable primary
comparison and misstated classical formulas; it is excluded
from the evidence.  Absent exact search hits, a bibliography,
and agreement of examples cannot establish originality.

\paragraph{Verdict submitted for independent review.}
C059PUBCOMPARISON is the mathematical comparison proved here,
conditional on the named published inputs and the retained
C059LOCAL proof.  C059PUBSIGNIFICANCE and
C059PUBCONTRIBUTION remain unestablished.  The Engineer
does not assign \texttt{novelty\_status=verified\_new},
\texttt{significance\_status=publishable}, or campaign completion.
The normal independent Reviewer should assess only these new
comparison and contribution claims and attribute the new
citations; the accepted proof and certificates need not be
reopened.  A found prior numerical evaluation would require
its exact reduction and replanning, not a weakened theorem or
an unrelated extension.

\paragraph{Source and computation fidelity.}
Primary PDFs and located excerpts are retained under
\texttt{research/literature/}.  The Compositio text layers omit
some displayed formulas; the formulas used here were read in
the rendered original pages, not reconstructed from search
summaries.  The previously confirmed Massey isolated-fibre
passage is reused from its existing literature archive.
The new source bindings await independent attribution; no
source is declared independently checked merely because its
identifier resolves.  The typed determinant output and exact
symbolic coefficient cancellations check only the indicated
algebra.  No Lean source or kernel evidence is introduced.

\begin{thebibliography}{9}
\bibitem{nemethi80}
A.~N\'emethi, \emph{Generalized local and global
Sebastiani--Thom type theorems}, Compositio Math.\ \textbf{80}
(1991), 1--14.  Theorem L1, Corollary 2.3, Examples 2.4,
and Section 3.2/Theorem G1.
\url{https://www.numdam.org/item/CM_1991__80_1_1_0/}.
\bibitem{nemethi-global}
A.~N\'emethi, \emph{Global Sebastiani--Thom theorem for polynomial
maps}, J.\ Math.\ Soc.\ Japan \textbf{43} (1991), 213--218.
Theorems 2.2 and 2.4.
\url{https://doi.org/10.2969/jmsj/04320213}.
\bibitem{nemethi79}
A.~N\'emethi, \emph{The Milnor fiber and the zeta function of the
singularities of type \(f=P(h,g)\)}, Compositio Math.\ \textbf{79}
(1991), 63--97.  Sections 2.0.1--2.0.4, Theorem A(a), and
Sections 2.1.4, 2.1.7.
\url{https://www.numdam.org/item/CM_1991__79_1_63_0/}.
\bibitem{le}
L\^e D\~ung Tr\'ang, \emph{Calcul du nombre de cycles
\'evanouissants d'une hypersurface complexe},
Ann.\ Inst.\ Fourier \textbf{23}(4) (1973), 261--270.
Theorems (1.1), (2.3), and Sections (3.1), (3.3).
\url{https://www.numdam.org/item/AIF_1973__23_4_261_0/}.
\bibitem{massey}
D.~B.~Massey, \emph{Non-isolated hypersurface singularities
and L\^e cycles}, arXiv:1410.3312v2 (November 13, 2015).
Section 1, items (4)--(5), Theorem 2.23, and Exercise 3.3.
\url{https://arxiv.org/abs/1410.3312v2}.
\bibitem{damon-pike}
J.~Damon and B.~Pike, \emph{Solvable groups, free divisors and
nonisolated matrix singularities II: Vanishing topology},
Geom.\ Topol.\ \textbf{18} (2014), 911--962.
Theorem 3.3, equation (3-2), p.~923.
\url{https://doi.org/10.2140/gt.2014.18.911}.
\end{thebibliography}
\end{document}
